\begin{solapartat}
\punts{0,35} L'energia del fotó en eV és:
\[ E_\textit{fotó} = 14,4\times 10^{3}\ \mathrm{eV}\,\frac{1,602\times 10^{-19}\ \mathrm{J}}{1\ \mathrm{eV}} = 2,31\times 10^{-15}\ \mathrm{J} \]

\punts{0,3} La freqüència del fotó és:
\[ f = \frac{E_\textit{fotó}}{h} = \frac{2,31\times 10^{-15}}{6,63\times 10^{-34}} = 3,48\times 10^{18}\ \mathrm{Hz} \]

\punts{0,3} La massa és:
\[ m = \frac{E_\textit{fotó}}{c^2} = \frac{2,31\times 10^{-15}}{(3,00\times 10^{8})^2} = 2,56\times 10^{-32}\ \mathrm{kg} \]

\punts{0,3} I la quantitat de moviment és:
\[ p = mc = 2,56\times 10^{-32}\cdot 3,00\times 10^{8} = 7,69\times 10^{-24}\ \mathrm{kg\,m\,s^{-1}} \]
\end{solapartat}

\begin{solapartat}
\punts{0,2} L'energia mecànica d'un fotó és la suma de l'energia cinètica i potencial gravitatòria:
\[ E_m = E_P + E_C = mgh + hf \]

\punts{0,4} Si imposem la conservació de l'energia, tenim que la diferència d'energia dels fotons és:
\[ E_{m,0} = E_{m,f} \;\Rightarrow\; mgh_0 + hf_0 = mgh_f + hf_f \]
\[ \Delta E_\textit{fotó} = hf_f - hf_0 = mgh_0 - mgh_f = 5,68\times 10^{-30}\ \mathrm{J} \]

Hem considerat com a punt inicial el punt elevat i, com a punt final, el punt més baix, per obtenir una variació positiva. En tot cas, el signe no és rellevant, atès que es demana el valor absolut de la variació d'energia i freqüència.

\punts{0,4} I la variació de la freqüència és:
\[ \Delta f = \frac{\Delta E_\textit{fotó}}{h} = \frac{5,68\times 10^{-30}}{6,63\times 10^{-34}} = 8,570\times 10^{3}\ \mathrm{Hz} \]

\punts{0,25} Atès que l'energia mecànica és la mateixa, i l'energia potencial gravitatòria és més gran en el punt elevat, l'energia del fotó i la seva freqüència seran més petites al punt més elevat. Per tant, el fotó té una freqüència més gran quan està a terra.
\end{solapartat}
